% ==================================================================
% Preâmbulo comum a todos os capítulos
% ==================================================================
\usepackage{fix-cm}
\usepackage[utf8]{inputenc}
\usepackage{babel}
\babelprovide[import,main]{brazilian}
\usepackage{amsmath,amssymb,amsthm,mathtools}
\usepackage[a4paper,margin=2.5cm]{geometry}
\usepackage{microtype}
\usepackage{anyfontsize}
\usepackage{enumitem}
\usepackage{booktabs,array}
\usepackage{xcolor}
\usepackage{pifont}
\usepackage{listings}
\usepackage{pgfplots}
\pgfplotsset{compat=1.16}
\usetikzlibrary{shapes.geometric,arrows.meta,positioning,fadings}
\usepackage[most]{tcolorbox}
\usepackage[colorlinks=true,linkcolor=azul,urlcolor=laranja,citecolor=azul,filecolor=laranja,
  bookmarksnumbered=true,pdftitle={Cálculo Numérico Interativo},pdfauthor={Osmar Alessio}]{hyperref}

\definecolor{azul}{RGB}{20,60,150}
\definecolor{azulescuro}{RGB}{10,30,80}
\definecolor{laranja}{RGB}{200,90,0}
\definecolor{verde}{RGB}{0,120,60}
\definecolor{vermelho}{RGB}{200,30,30}
\definecolor{cL1}{RGB}{31,99,180}
\definecolor{cL2}{RGB}{27,138,59}
\definecolor{cL3}{RGB}{230,126,34}
\definecolor{cL4}{RGB}{148,103,189}

\sloppy
\setlength{\parskip}{0.4em}
\renewcommand{\arraystretch}{1.15}

\lstset{basicstyle=\ttfamily\small,keywordstyle=\color{azul}\bfseries,commentstyle=\color{gray},
  stringstyle=\color{verde},frame=single,rulecolor=\color{gray!50},backgroundcolor=\color{gray!4},
  columns=fullflexible,keepspaces=true,showstringspaces=false,breaklines=true}

% ------------------------------------------------------------------
% Links para os aplicativos do GeoGebra e para os textos complementares.
% Opção 1 (padrão): arquivos .ggb na pasta ggb/ e PDFs na pasta complementos/,
% ao lado do PDF do livro.
% Opção 2: depois de publicar os arquivos no geogebra.org, troque
% \ggbabrir por \ggbonline e informe o endereço de cada material.
% Os nomes dos arquivos de cada capítulo (\arq...) são definidos no início
% de cada arquivo da pasta capitulos/.
% ------------------------------------------------------------------
\newcommand{\ggbpasta}{ggb/}
\newcommand{\ggbabrir}[2]{\href{run:\ggbpasta #1}{#2}}
\newcommand{\ggbonline}[2]{\href{#1}{#2}}
\newcommand{\pdfabrir}[2]{\href{complementos/#1}{#2}}
\newcommand{\arqtxt}[1]{\texttt{\expandafter\detokenize\expandafter{#1}}}

\newtcolorbox{atividade}[2]{enhanced,breakable,colback=orange!6,colframe=laranja,
  fonttitle=\bfseries,coltitle=white,
  title={Atividade no GeoGebra: #1\hfill\ggbabrir{#2}{\textcolor{white}{[abrir o aplicativo]}}},
  after upper={\par\smallskip\small\textbf{Arquivo:}\\ \ggbabrir{#2}{\texttt{ggb/\arqtxt{#2}}}}}

\newtcolorbox{leitura}[2]{enhanced,breakable,colback=teal!4,colframe=teal!70!black,
  fonttitle=\bfseries,coltitle=white,
  title={Leitura complementar: #1\hfill\pdfabrir{#2}{\textcolor{white}{[abrir o texto]}}},
  after upper={\par\smallskip\small\textbf{Arquivo:} \pdfabrir{#2}{\texttt{complementos/\arqtxt{#2}}}}}

\newtcolorbox{exemplo}[1][]{enhanced,breakable,colback=blue!3,colframe=azul!70,
  fonttitle=\bfseries,title={Exemplo#1}}

\newtcolorbox{observacao}{enhanced,breakable,colback=green!4,colframe=verde!70,
  fonttitle=\bfseries,title={Observação}}

\theoremstyle{definition}
\newtheorem{definicao}{Definição}[chapter]
\theoremstyle{plain}
\newtheorem{teorema}[definicao]{Teorema}
\newtheorem{proposicao}[definicao]{Proposição}
\newtheorem{corolario}[definicao]{Corolário}

\tcolorboxenvironment{teorema}{colback=azul!4,colframe=azul,boxrule=0.6pt,arc=1mm,breakable}
\tcolorboxenvironment{proposicao}{colback=azul!4,colframe=azul,boxrule=0.6pt,arc=1mm,breakable}
\tcolorboxenvironment{definicao}{colback=gray!5,colframe=gray!60,boxrule=0.6pt,arc=1mm,breakable}

\renewcommand{\qedsymbol}{$\blacksquare$}
\newcommand{\x}{\mathbf{x}}
\newcommand{\bb}{\mathbf{b}}
\newcommand{\y}{\mathbf{y}}

% Dados do livro (edite aqui)
\newcommand{\livrotitulo}{Cálculo Numérico Interativo}
\newcommand{\livrosubtitulo}{Aritmética de ponto flutuante, zeros de funções,\\sistemas lineares, interpolação, ajuste de curvas\\e integração numérica com aplicativos no GeoGebra}
\newcommand{\livroautor}{Osmar Alessio}
\newcommand{\livroinstituicao}{Universidade Federal do Triângulo Mineiro}
\newcommand{\livrocurso}{Licenciatura em Matemática -- ICENE}
\newcommand{\livrolocal}{Uberaba -- MG}
\newcommand{\livroano}{2026}
